\par\noindent\begin{minipage}{\linewidth}
% !TeX root = ../教材题库_全部条目_转写工作稿.tex
% Source: XJ-B1-C02-HW-P049-Y0879
% 第2章 / 2.2 从函数观点看一元二次方程 / 教材 p.49 / PDF 第56页
% 已对照原页转写题干及公式；待教师复核；未制作答案
% 原题及相邻原解截图（核对用）：../题目原图/XJ-B1-C02-HW-P049-Y0879-01.webp
% 本文件只在转写完成后作为题干使用；OCR 条目尚不是可靠题干。
\begin{bankitem}{习题2.2 第5题}{第2章 / 2.2 从函数观点看一元二次方程 / 教材 p.49 / PDF 第56页}{XJ-B1-C02-HW-P049-Y0879}{已对照原页转写题干及公式；待教师复核；未制作答案}
已知二次函数 $y=(1-m)x^2+4x-3$ 的图象开口向下，与 $x$ 轴交于 $A(x_1,0),B(x_2,0)$ 两点。
\subq{1}求 $m$ 的取值范围；
\subq{2}当 $x_1^2+x_2^2=10$ 时，求该二次函数的表达式。
\end{bankitem}
\end{minipage}\par\medskip
